\documentclass[11pt]{article}

\usepackage[margin=1in]{geometry}
\usepackage[T1]{fontenc}
\usepackage[utf8]{inputenc}
\usepackage{lmodern}
\usepackage{amsmath,amssymb,amsthm,mathtools,bm}
\usepackage{booktabs}
\usepackage{array}
\usepackage{hyperref}
\usepackage{enumitem}
\usepackage{xcolor}
\usepackage{listings}
\usepackage{float}
\usepackage{microtype}

\hypersetup{
  colorlinks=true,
  linkcolor=blue,
  urlcolor=blue,
  citecolor=blue,
  pdftitle={Monitoring Diagonal Spectral Mass and Structure Factors in a Transformer for Fibonacci Modulo n},
  pdfauthor={OpenAI}
}

\definecolor{codebg}{RGB}{247,247,247}
\definecolor{codecomment}{RGB}{0,110,40}
\definecolor{codekeyword}{RGB}{0,0,150}
\definecolor{codestring}{RGB}{160,20,20}

\lstdefinestyle{pythonstyle}{
  language=Python,
  basicstyle=\ttfamily\small,
  keywordstyle=\color{codekeyword}\bfseries,
  commentstyle=\color{codecomment},
  stringstyle=\color{codestring},
  backgroundcolor=\color{codebg},
  showstringspaces=false,
  breaklines=true,
  frame=single,
  rulecolor=\color{black!20},
  tabsize=4,
  columns=fullflexible,
  keepspaces=true
}

\newcommand{\Zn}{\mathbb{Z}_n}
\newcommand{\E}{\mathbb{E}}
\newcommand{\R}{\mathbb{R}}
\newcommand{\C}{\mathbb{C}}
\newcommand{\ii}{\mathrm{i}}
\newcommand{\diag}{\mathrm{diag}}
\newcommand{\Tr}{\mathrm{Tr}}
\newcommand{\Var}{\mathrm{Var}}
\newcommand{\Cov}{\mathrm{Cov}}

\title{Monitoring Diagonal Spectral Mass and Structure Factors During Training\\[4pt]
\large A student-friendly guide for a Transformer trained on Fibonacci modulo $n$}
\author{Prepared for the Phase Transitions in Neural Networks project}
\date{March 2026}

\begin{document}
\maketitle

\begin{abstract}
This note explains, in mathematical and conceptual detail, two quantities that can be monitored during training of a Transformer that predicts the Fibonacci recurrence modulo $n$:
\emph{diagonal spectral mass} and the \emph{structure factor}. The main goals are:
\begin{enumerate}[label=(\roman*)]
    \item to define these quantities carefully in a way that is natural for the task,
    \item to explain why they are related but not identical,
    \item to show how the Ornstein-Zernike form of the structure factor appears near a continuous phase transition,
    \item to provide ready-to-use Python code, and
    \item to explain how to run a scaling study by training the model for different values of $n$.
\end{enumerate}
The intended audience is college students with basic calculus. For that reason, the explanations are detailed and avoid unnecessary abstraction.
\end{abstract}

\tableofcontents
\newpage

\section{The learning problem and the right state space}

We consider sequences generated by the Fibonacci recurrence modulo $n$:
\begin{equation}
    x_{t+1} = x_t + x_{t-1} \pmod n.
\end{equation}
Each token belongs to the vocabulary
\begin{equation}
    \Zn = \{0,1,\dots,n-1\}.
\end{equation}

A Transformer is trained autoregressively: given a prefix $x_1,x_2,\dots,x_t$, it predicts the next token $x_{t+1}$.

\subsection{Why the pair $(x_{t-1},x_t)$ matters}

For this task, the correct next token depends only on the last two tokens, not on the whole prefix. The sufficient state is
\begin{equation}
    s_t = (x_{t-1},x_t) \in \Zn^2.
\end{equation}
Once we know $s_t=(a,b)$, the correct target is
\begin{equation}
    y = a+b \pmod n.
\end{equation}

So the natural ``space of states'' is not the one-dimensional vocabulary $\Zn$, but the two-dimensional discrete torus
\begin{equation}
    \Zn^2 = \{(a,b): a,b \in \Zn\}.
\end{equation}
It has $n^2$ points.

This point is essential. If we want to study how the network learns the rule, then the most informative object is usually a field defined on this pair-state lattice. That is the setting in which diagonal spectral mass and the structure factor become natural.

\subsection{A canonical probe set}

To evaluate the model at a fixed training step, we would like to know what it predicts for every possible pair $(a,b)$. Because the recurrence is deterministic, we can construct a valid prefix ending in $(a,b)$ by running the recurrence backward:
\begin{equation}
    x_{t-2} = x_t - x_{t-1} \pmod n.
\end{equation}
Repeating this allows us to build one valid probe sequence for each of the $n^2$ states.

So for every checkpoint in training, we can build a table of model outputs indexed by $(a,b)$.

\section{From logits to a field on the pair lattice}

Suppose the model outputs logits
\begin{equation}
    \ell_\theta(c \mid a,b), \qquad c \in \Zn,
\end{equation}
when fed a canonical valid prefix ending in $(a,b)$. Here $\theta$ denotes the model parameters at a given training step.

For fixed $c$, the quantity
\begin{equation}
    \phi_c(a,b) = \ell_\theta(c \mid a,b)
\end{equation}
is a scalar field on the discrete torus $\Zn^2$.

\subsection{Gauge issue: why logits should be centered first}

Raw logits have an important redundancy. If we add the same number $g(a,b)$ to all class logits at the state $(a,b)$,
\begin{equation}
    \ell_\theta(c \mid a,b) \mapsto \ell_\theta(c \mid a,b) + g(a,b),
\end{equation}
then the softmax probabilities do not change. Therefore, some changes in the raw logit field are not physically meaningful for prediction.

To remove this redundancy, we define centered logits:
\begin{equation}
    \widetilde{\ell}_\theta(c \mid a,b)
    = \ell_\theta(c \mid a,b) - \frac{1}{n}\sum_{c'=0}^{n-1} \ell_\theta(c' \mid a,b).
\end{equation}
This subtracts the mean across classes at each pair $(a,b)$.

After this centering step, the field reflects only relative class preferences.

\subsection{Other possible fields}

The centered logits are not the only possible field. You may also define a vector field from a hidden representation,
\begin{equation}
    h_\theta(a,b) \in \R^d,
\end{equation}
for example the final hidden state at the last position of the sequence ending in $(a,b)$.

Then one can study:
\begin{itemize}
    \item a scalar projection $u^\top h_\theta(a,b)$ for a fixed vector $u$, or
    \item the rotationally invariant correlation function
    \begin{equation}
        C(r) = \frac{1}{n^2} \sum_{s \in \Zn^2} \widetilde{h}(s) \cdot \widetilde{h}(s+r),
    \end{equation}
    where $\widetilde{h}$ is centered over the lattice.
\end{itemize}

For the first implementation, centered logits are often the easiest choice because they are already produced by the model and are directly tied to prediction.

\section{Discrete Fourier analysis on $\Zn^2$}

\subsection{Why Fourier modes are useful}

If the model learns a spatial pattern on the pair lattice, Fourier analysis tells us \emph{which frequencies} carry most of the signal. This is directly analogous to physics, where one often studies ordering by asking whether power condenses near particular wavevectors.

A two-dimensional discrete Fourier transform is the natural tool because our state space is the finite grid $\Zn^2$.

\subsection{The 2D discrete Fourier transform}

Let $\phi(a,b)$ be a scalar field on $\Zn^2$. Define
\begin{equation}
    \omega = e^{2\pi \ii / n}.
\end{equation}
The discrete Fourier transform is
\begin{equation}
    \widehat{\phi}(k,\ell)
    = \frac{1}{n} \sum_{a=0}^{n-1} \sum_{b=0}^{n-1} \phi(a,b) \, \omega^{-ka-\ell b},
    \qquad k,\ell \in \Zn.
\end{equation}
The inverse transform is
\begin{equation}
    \phi(a,b)
    = \frac{1}{n} \sum_{k=0}^{n-1} \sum_{\ell=0}^{n-1} \widehat{\phi}(k,\ell) \, \omega^{ka+\ell b}.
\end{equation}

The coefficients are complex numbers, but this is not a problem. In practice, the physically relevant quantity is often the \emph{power}
\begin{equation}
    |\widehat{\phi}(k,\ell)|^2
    = \widehat{\phi}(k,\ell) \overline{\widehat{\phi}(k,\ell)}
    = (\Re \widehat{\phi})^2 + (\Im \widehat{\phi})^2.
\end{equation}

So although the Fourier coefficients themselves are complex, the power spectrum is real and nonnegative.

\subsection{Fourier transform of the centered logit table}

Since the logits depend on the output class $c$, we can Fourier transform each class slice separately:
\begin{equation}
    \widehat{\widetilde{\ell}}_\theta(k,\ell;c)
    = \frac{1}{n} \sum_{a=0}^{n-1}\sum_{b=0}^{n-1}
    \widetilde{\ell}_\theta(c \mid a,b)\, \omega^{-ka-\ell b}.
\end{equation}
The corresponding class-resolved power is
\begin{equation}
    P_\theta(k,\ell;c) = \left| \widehat{\widetilde{\ell}}_\theta(k,\ell;c) \right|^2.
\end{equation}
Summing over $c$ produces a single nonnegative 2D spectrum:
\begin{equation}
    S_\theta(k,\ell) = \sum_{c=0}^{n-1} P_\theta(k,\ell;c).
\end{equation}
This $S_\theta$ will be our first version of the structure factor.

\section{Why diagonal modes are special for Fibonacci modulo $n$}

The correct label at state $(a,b)$ is $c=a+b \pmod n$. The indicator of the correct class can be written exactly as
\begin{equation}
    \mathbf{1}\{c = a+b \pmod n\}
    = \frac{1}{n} \sum_{m=0}^{n-1} \omega^{m(a+b-c)}.
\end{equation}
Now look only at the dependence on $(a,b)$:
\begin{equation}
    \omega^{m(a+b-c)} = \omega^{-mc} \omega^{ma} \omega^{mb}.
\end{equation}
The Fourier dependence on the two lattice directions uses the \emph{same} mode number $m$ in both directions. Therefore the exact target rule lives on the diagonal set
\begin{equation}
    (k,\ell) = (m,m).
\end{equation}
This is the mathematical reason that diagonal spectral mass is a natural quantity for this task.

\section{Definition of diagonal spectral mass}

Let
\begin{equation}
    S_\theta(k,\ell) = \sum_{c=0}^{n-1} \left|\widehat{\widetilde{\ell}}_\theta(k,\ell;c)\right|^2.
\end{equation}
The zero-frequency term $(0,0)$ is called the \emph{DC mode}. After centering logits across classes, it is much less problematic, but in spectral order parameters it is standard to exclude the DC contribution because it measures a uniform background rather than nontrivial spatial organization.

We define the diagonal spectral mass by
\begin{equation}
    M_{\diag}(\theta)
    = \frac{\sum_{m=1}^{n-1} S_\theta(m,m)}{\sum_{(k,\ell) \neq (0,0)} S_\theta(k,\ell)}.
\end{equation}
This is a dimensionless number between $0$ and $1$.

\subsection{Interpretation}

\begin{itemize}
    \item If the spectrum is spread out over many unrelated modes, then $M_{\diag}$ is small.
    \item If most of the power sits on diagonal modes $(m,m)$, then $M_{\diag}$ is large.
    \item If the network learns the modular addition structure cleanly, one expects $M_{\diag}$ to increase during training.
\end{itemize}

\subsection{What it does and does not measure}

This quantity tells us \emph{where} the spectral weight sits, but not how sharply it is concentrated around a single mode. That is why it is order-parameter-like, but not itself a correlation length.

\section{Definition of the structure factor}

\subsection{Two-point correlation function}

Suppose we have a scalar field $\phi(s)$ defined on lattice points $s \in \Zn^2$. The connected two-point correlation function is
\begin{equation}
    C(r) = \frac{1}{n^2}\sum_{s \in \Zn^2} \left(\phi(s)-\bar{\phi}\right)\left(\phi(s+r)-\bar{\phi}\right),
\end{equation}
where
\begin{equation}
    \bar{\phi} = \frac{1}{n^2}\sum_{s \in \Zn^2} \phi(s).
\end{equation}
The vector $r$ is a displacement on the torus.

Conceptually, $C(r)$ measures how similar the field is to a shifted copy of itself. If $C(r)$ stays large for large $r$, then correlations are long-ranged. If it falls quickly, then correlations are short-ranged.

\subsection{Fourier transform of the correlation function}

The structure factor is the Fourier transform of the correlation function:
\begin{equation}
    S(k) = \sum_{r \in \Zn^2} C(r)e^{-\ii k \cdot r}.
\end{equation}
In the discrete setting, we can index $k$ by pairs $(k,\ell) \in \Zn^2$.

There is an equivalent and very useful identity:
\begin{equation}
    S(k) = \left|\widehat{\phi}(k)\right|^2
\end{equation}
when $\phi$ has been centered so that its mean is zero. More precisely, up to normalization conventions,
\begin{equation}
    S(k) \propto \left|\widehat{\phi}(k)\right|^2.
\end{equation}
This is a discrete version of the Wiener-Khinchin theorem.

So the structure factor can be viewed in two equivalent ways:
\begin{enumerate}
    \item Fourier transform of the correlation function.
    \item Power spectrum of the centered field.
\end{enumerate}

\subsection{Structure factor from logits}

For the centered logit field, a natural definition is
\begin{equation}
    S_\theta(k,\ell) = \sum_{c=0}^{n-1} \left|\widehat{\widetilde{\ell}}_\theta(k,\ell;c)\right|^2.
\end{equation}
This combines the contribution from all output classes into one spectrum on the pair lattice.

\subsection{Structure factor from hidden states}

If we instead use a vector field $h(s) \in \R^d$, then a natural rotationally invariant structure factor is
\begin{equation}
    S_h(k) = \sum_{j=1}^{d} \left| \widehat{\widetilde{h}_j}(k) \right|^2,
\end{equation}
where $\widetilde{h}_j$ is the centered $j$-th coordinate of the hidden state. This is often a very good choice when we want to study learned internal representations rather than only the final prediction layer.

\section{How diagonal spectral mass and structure factor are related}

The two quantities are closely related, but they answer different questions.

\subsection{Diagonal mass is a summary statistic built from the structure factor}

Once the full structure factor $S_\theta(k,\ell)$ is known, the diagonal spectral mass is obtained by summing its weight on the diagonal and dividing by the total nonzero spectral weight:
\begin{equation}
    M_{\diag}(\theta)
    = \frac{\sum_{m=1}^{n-1} S_\theta(m,m)}{\sum_{(k,\ell)\neq(0,0)} S_\theta(k,\ell)}.
\end{equation}
So diagonal spectral mass is not an independent object. It is a particular scalar functional of the full structure factor.

\subsection{What the full structure factor contains that diagonal mass loses}

The structure factor retains the entire two-dimensional shape of the spectrum:
\begin{itemize}
    \item where the dominant peak is,
    \item how wide it is,
    \item whether multiple peaks coexist,
    \item whether the spectrum is isotropic or elongated.
\end{itemize}
By contrast, $M_{\diag}$ collapses all of that information to a single number.

For example, two checkpoints may have the same diagonal spectral mass, but one could have:
\begin{enumerate}
    \item broad diffuse power spread along the whole diagonal, while the other has
    \item a very sharp peak at one diagonal mode.
\end{enumerate}
These are physically different situations. The correlation length depends strongly on the width of the peak, not only on the total diagonal weight.

\subsection{Practical consequence}

It is wise to monitor both:
\begin{itemize}
    \item $M_{\diag}(\theta)$ as a simple order-parameter-like summary, and
    \item $S_\theta(k,\ell)$ as the full spectral diagnostic from which one can estimate correlation lengths.
\end{itemize}

\section{Correlation length and why it matters}

In many continuous phase transitions, the key idea is that local pieces of the system become correlated over larger and larger distances. The characteristic scale over which this happens is the \emph{correlation length}, usually denoted by $\xi$.

If the correlation function behaves approximately like
\begin{equation}
    C(r) \sim e^{-|r|/\xi},
\end{equation}
then:
\begin{itemize}
    \item small $\xi$ means short-range order,
    \item large $\xi$ means long-range coherence,
    \item $\xi \to \infty$ is the idealized critical behavior in an infinite system.
\end{itemize}

In a finite system, $\xi$ cannot literally diverge to infinity, but it can become comparable to the system size. In our setup, the natural linear size is $n$, because the state space is an $n \times n$ torus.

\section{Where the Ornstein-Zernike form comes from}

This section gives an intuitive calculus-level derivation of the standard formula
\begin{equation}
    S(k) \approx \frac{A}{\xi^{-2} + |k-k^*|^2}
\end{equation}
near a continuous phase transition.

\subsection{Step 1: expect a dominant ordering wavevector}

In many systems the structure factor develops a peak around some preferred wavevector $k^*$. In our Fibonacci setting, the preferred modes are expected to lie on the diagonal, and eventually perhaps near a particular diagonal wavevector or a small set of them.

So instead of studying $S(k)$ everywhere, we zoom in near the dominant peak. Write
\begin{equation}
    q = k-k^*.
\end{equation}
Then the question becomes: what does the structure factor look like for small $q$?

\subsection{Step 2: short summary of the physical idea}

Near a continuous transition, the field varies slowly on large scales. For slowly varying fields, the leading spatial term in an energy functional is quadratic in gradients. In real space that means a typical linear response equation has the form
\begin{equation}
    \left(-\nabla^2 + \xi^{-2}\right) C(r) = A\,\delta(r).
\end{equation}
Here:
\begin{itemize}
    \item $-\nabla^2$ penalizes rapid spatial variation,
    \item $\xi^{-2}$ sets the inverse square correlation length,
    \item $\delta(r)$ is a point source,
    \item $A$ is a constant amplitude.
\end{itemize}

This equation is the simplest rotationally symmetric linear equation one can write down for long-wavelength fluctuations. It is the continuum version of a lattice theory with nearest-neighbor coupling.

\subsection{Step 3: Fourier transform the equation}

The Fourier transform of $-\nabla^2$ is $|k|^2$. Therefore, taking the Fourier transform of
\begin{equation}
    \left(-\nabla^2 + \xi^{-2}\right) C(r) = A\,\delta(r)
\end{equation}
produces
\begin{equation}
    \left(|k|^2 + \xi^{-2}\right) \widehat{C}(k) = A.
\end{equation}
So
\begin{equation}
    \widehat{C}(k) = \frac{A}{\xi^{-2}+|k|^2}.
\end{equation}
But the structure factor is the Fourier transform of the correlation function, so
\begin{equation}
    S(k) \approx \frac{A}{\xi^{-2}+|k|^2}.
\end{equation}
If the dominant ordering occurs around $k^*$ rather than around $0$, we replace $k$ by $k-k^*$:
\begin{equation}
    S(k) \approx \frac{A}{\xi^{-2}+|k-k^*|^2}.
\end{equation}
This is the Ornstein-Zernike form.

\subsection{Step 4: geometric interpretation}

The denominator tells us that the peak width is controlled by $\xi^{-1}$:
\begin{itemize}
    \item when $\xi$ is small, the peak is broad,
    \item when $\xi$ is large, the peak is narrow and tall.
\end{itemize}
So a growing correlation length means the spectrum becomes increasingly concentrated near the dominant mode.

\subsection{Step 5: relation with exponential decay in real space}

One can also understand this from the opposite direction. If real-space correlations decay exponentially over a length scale $\xi$, their Fourier transform is a smooth peak whose width is roughly of order $1/\xi$. That is why measuring the peak width of the structure factor is a standard way to estimate correlation length.

\subsection{Important caution}

The Ornstein-Zernike form is an approximation valid near a continuous transition and for long wavelengths. It does not mean the exact finite-size lattice spectrum has that form at every point. In practice, we use it locally, near the dominant peak, to motivate a numerical estimator of $\xi$.

\section{Second-moment estimator of the correlation length}

Because our system lives on a finite periodic lattice, we estimate $\xi$ from the structure factor using nearby lattice momenta.

\subsection{Lattice momenta}

The allowed momenta are
\begin{equation}
    k_x = \frac{2\pi m_x}{n}, \qquad k_y = \frac{2\pi m_y}{n},
\end{equation}
with $m_x,m_y \in \Zn$.

Suppose the dominant nonzero peak occurs at $k^*$. Let $\delta k$ denote one nearest-neighbor step in momentum space, for example
\begin{equation}
    \delta k_x = \left(\frac{2\pi}{n},0\right), \qquad
    \delta k_y = \left(0,\frac{2\pi}{n}\right).
\end{equation}
Using the Ornstein-Zernike form,
\begin{equation}
    S(k) \approx \frac{A}{\xi^{-2}+|k-k^*|^2},
\end{equation}
we get approximately
\begin{equation}
    S(k^*) \approx A\xi^2,
\end{equation}
\begin{equation}
    S(k^*+\delta k) \approx \frac{A}{\xi^{-2}+|\delta k|^2}.
\end{equation}
Taking the ratio gives
\begin{equation}
    \frac{S(k^*)}{S(k^*+\delta k)} \approx 1 + \xi^2 |\delta k|^2.
\end{equation}
Therefore,
\begin{equation}
    \xi^2 \approx \frac{1}{|\delta k|^2}\left(\frac{S(k^*)}{S(k^*+\delta k)} - 1\right).
\end{equation}
On a periodic lattice it is more accurate to replace $|\delta k|^2$ by the lattice expression
\begin{equation}
    4\sin^2\left(\frac{\pi}{n}\right).
\end{equation}
Thus a common estimator is
\begin{equation}
    \xi^2 \approx \frac{1}{4\sin^2(\pi/n)}
    \left(\frac{S(k^*)}{\overline{S}_{\text{nn}}(k^*)} - 1\right),
\end{equation}
where $\overline{S}_{\text{nn}}(k^*)$ is the average of the structure factor over the nearest momentum neighbors of $k^*$.

\section{A practical monitoring strategy during training}

At each saved checkpoint in training, perform the following steps:
\begin{enumerate}
    \item Build the $n^2$ canonical probe sequences, one for each pair $(a,b)$.
    \item Run the model on these sequences.
    \item Extract the final-position logits and reshape them into a tensor
    \begin{equation}
        L[a,b,c].
    \end{equation}
    \item Center the logits across classes to obtain $\widetilde{L}[a,b,c]$.
    \item Apply a 2D FFT over the $(a,b)$ axes.
    \item Compute the class-summed structure factor
    \begin{equation}
        S[k,\ell] = \sum_c \left|\widehat{\widetilde{L}}[k,\ell,c]\right|^2.
    \end{equation}
    \item Compute the diagonal spectral mass from $S$.
    \item Locate the dominant nonzero mode and estimate the correlation length from the peak width.
\end{enumerate}

If you do this repeatedly during training, you obtain trajectories such as
\begin{equation}
    M_{\diag}(t), \qquad \xi(t), \qquad S_t(k,\ell),
\end{equation}
where $t$ is training step or epoch.

\section{Python implementation}

The code below uses PyTorch. It assumes that calling the model on an integer tensor of shape \texttt{[batch, seq\_len]} returns logits of shape \texttt{[batch, seq\_len, n]}.

\subsection{Building the canonical probe sequences}

\begin{lstlisting}[style=pythonstyle]
import math
import torch


def build_probe_sequences(n: int, seq_len: int, device=None):
    """
    Construct one valid Fibonacci-mod-n sequence for each terminal pair (a, b).

    Returns
    -------
    seqs : LongTensor of shape [n*n, seq_len]
        Canonical valid sequences ending in pair (a, b).
    pairs : LongTensor of shape [n*n, 2]
        The final pair for each sequence.
    """
    vals = torch.arange(n, device=device)
    pairs = torch.cartesian_prod(vals, vals)   # [n*n, 2]

    seqs = torch.zeros((n * n, seq_len), dtype=torch.long, device=device)
    seqs[:, -2] = pairs[:, 0]
    seqs[:, -1] = pairs[:, 1]

    # Run the recurrence backward:
    # x_i = x_{i+2} - x_{i+1} (mod n)
    for i in range(seq_len - 3, -1, -1):
        seqs[:, i] = (seqs[:, i + 2] - seqs[:, i + 1]) % n

    return seqs, pairs
\end{lstlisting}

\subsection{Collecting the logit table $L[a,b,c]$}

\begin{lstlisting}[style=pythonstyle]
@torch.no_grad()
def collect_state_logits(model, n: int, seq_len: int, device, batch_size: int = 256):
    """
    Evaluate the model on every pair-state (a, b) and collect the final logits.

    Returns
    -------
    L : FloatTensor of shape [n, n, n]
        L[a, b, c] is the logit for class c when the sequence ends in (a, b).
    """
    model.eval()
    seqs, pairs = build_probe_sequences(n=n, seq_len=seq_len, device=device)
    L = torch.empty((n, n, n), dtype=torch.float32, device=device)

    for start in range(0, seqs.shape[0], batch_size):
        x = seqs[start:start + batch_size]
        logits = model(x)                  # expected shape [B, seq_len, n]
        final_logits = logits[:, -1, :]    # shape [B, n]

        ab = pairs[start:start + batch_size]
        a = ab[:, 0]
        b = ab[:, 1]
        L[a, b, :] = final_logits

    return L
\end{lstlisting}

\subsection{Centering logits and computing the full structure factor}

\begin{lstlisting}[style=pythonstyle]
def structure_factor_from_logits_table(L: torch.Tensor):
    """
    Compute the class-summed structure factor from a logit table.

    Parameters
    ----------
    L : FloatTensor of shape [n, n, n]
        Axes are [a, b, c].

    Returns
    -------
    S : FloatTensor of shape [n, n]
        Class-summed power spectrum on the pair-state lattice.
    L_centered : FloatTensor of shape [n, n, n]
        Centered logits.
    F : ComplexTensor of shape [n, n, n]
        2D FFT of the centered logits over the (a, b) axes.
    """
    n = L.shape[0]
    assert L.shape == (n, n, n)

    # Remove the softmax gauge freedom.
    L_centered = L - L.mean(dim=-1, keepdim=True)

    # 2D FFT over pair-state axes only.
    F = torch.fft.fftn(L_centered, dim=(0, 1), norm="ortho")

    # Class-resolved power spectrum.
    P = F.abs().pow(2)

    # Sum over the class index c.
    S = P.sum(dim=-1)
    return S, L_centered, F
\end{lstlisting}

\subsection{Computing the diagonal spectral mass}

\begin{lstlisting}[style=pythonstyle]
def diagonal_spectral_mass_from_structure_factor(S: torch.Tensor):
    """
    Compute diagonal spectral mass from a 2D structure factor S[k, l].
    """
    n = S.shape[0]
    assert S.shape == (n, n)

    idx = torch.arange(n, device=S.device)

    dc = S[0, 0]
    total_nonzero = S.sum() - dc
    diag_nonzero = S[idx, idx].sum() - dc

    return diag_nonzero / (total_nonzero + 1e-12)


def diagonal_spectral_mass_from_logits_table(L: torch.Tensor):
    S, _, _ = structure_factor_from_logits_table(L)
    return diagonal_spectral_mass_from_structure_factor(S)
\end{lstlisting}

\subsection{Locating the dominant nonzero mode}

\begin{lstlisting}[style=pythonstyle]
def dominant_nonzero_mode(S: torch.Tensor):
    """
    Return the lattice indices (kx_idx, ky_idx) of the largest nonzero mode.
    """
    n = S.shape[0]
    S_work = S.clone()
    S_work[0, 0] = -torch.inf

    flat_idx = torch.argmax(S_work)
    kx_idx = int(flat_idx // n)
    ky_idx = int(flat_idx % n)
    return kx_idx, ky_idx
\end{lstlisting}

\subsection{Estimating the correlation length from the peak width}

\begin{lstlisting}[style=pythonstyle]
def second_moment_correlation_length(S: torch.Tensor):
    """
    Estimate the correlation length from the dominant nonzero peak of S.

    Uses the nearest momentum neighbors of the dominant mode and a lattice
    version of the Ornstein-Zernike second-moment estimator.
    """
    n = S.shape[0]
    kx, ky = dominant_nonzero_mode(S)

    S_peak = S[kx, ky]

    neighbors = [
        ((kx + 1) % n, ky),
        ((kx - 1) % n, ky),
        (kx, (ky + 1) % n),
        (kx, (ky - 1) % n),
    ]

    S_nn = torch.stack([S[i, j] for i, j in neighbors]).mean()

    denom = 4.0 * math.sin(math.pi / n) ** 2
    xi_sq = (S_peak / (S_nn + 1e-12) - 1.0) / (denom + 1e-12)
    xi_sq = torch.clamp(xi_sq, min=0.0)
    xi = torch.sqrt(xi_sq)

    return {
        "xi": xi,
        "peak_index": (kx, ky),
        "S_peak": S_peak,
        "S_nn_mean": S_nn,
    }
\end{lstlisting}

\subsection{A single monitoring function}

\begin{lstlisting}[style=pythonstyle]
@torch.no_grad()
def monitor_spectral_observables(model, n: int, seq_len: int, device, batch_size: int = 256):
    L = collect_state_logits(model, n=n, seq_len=seq_len, device=device, batch_size=batch_size)
    S, _, _ = structure_factor_from_logits_table(L)
    m_diag = diagonal_spectral_mass_from_structure_factor(S)
    xi_info = second_moment_correlation_length(S)

    return {
        "structure_factor": S,
        "diagonal_spectral_mass": m_diag.item(),
        "correlation_length": xi_info["xi"].item(),
        "peak_index": xi_info["peak_index"],
        "S_peak": xi_info["S_peak"].item(),
        "S_nn_mean": xi_info["S_nn_mean"].item(),
    }
\end{lstlisting}

\subsection{How to integrate the monitor into a training loop}

\begin{lstlisting}[style=pythonstyle]
monitor_every = 100
history = {
    "step": [],
    "diag_mass": [],
    "xi": [],
    "peak_kx": [],
    "peak_ky": [],
}

for step in range(num_training_steps):
    model.train()
    optimizer.zero_grad()

    logits = model(x_batch)
    loss = loss_fn(logits[:, :-1, :].reshape(-1, n), y_batch[:, :-1].reshape(-1))
    loss.backward()
    optimizer.step()

    if step % monitor_every == 0:
        obs = monitor_spectral_observables(
            model=model,
            n=n,
            seq_len=seq_len,
            device=device,
            batch_size=256,
        )

        history["step"].append(step)
        history["diag_mass"].append(obs["diagonal_spectral_mass"])
        history["xi"].append(obs["correlation_length"])
        history["peak_kx"].append(obs["peak_index"][0])
        history["peak_ky"].append(obs["peak_index"][1])

        print(
            f"step={step:6d}  "
            f"diag_mass={obs['diagonal_spectral_mass']:.6f}  "
            f"xi={obs['correlation_length']:.4f}  "
            f"peak={obs['peak_index']}"
        )
\end{lstlisting}

\subsection{Plotting the observables}

\begin{lstlisting}[style=pythonstyle]
import matplotlib.pyplot as plt


def plot_history(history):
    steps = history["step"]

    plt.figure(figsize=(7, 4))
    plt.plot(steps, history["diag_mass"], marker="o")
    plt.xlabel("training step")
    plt.ylabel("diagonal spectral mass")
    plt.title("Diagonal spectral mass during training")
    plt.tight_layout()
    plt.show()

    plt.figure(figsize=(7, 4))
    plt.plot(steps, history["xi"], marker="o")
    plt.xlabel("training step")
    plt.ylabel("correlation length estimator")
    plt.title("Estimated correlation length during training")
    plt.tight_layout()
    plt.show()
\end{lstlisting}

\subsection{Visualizing the structure factor itself}

\begin{lstlisting}[style=pythonstyle]
def plot_structure_factor(S: torch.Tensor, title: str = "Structure factor"):
    plt.figure(figsize=(5, 4))
    plt.imshow(S.detach().cpu().numpy(), origin="lower", aspect="auto")
    plt.colorbar(label="power")
    plt.xlabel("k_y index")
    plt.ylabel("k_x index")
    plt.title(title)
    plt.tight_layout()
    plt.show()
\end{lstlisting}

Seeing the full heatmap is extremely useful. A rising diagonal spectral mass may come from a broad diagonal band, while a growing correlation length should show a sharper and sharper localized peak.

\section{How to interpret the monitored quantities}

\subsection{If $M_{\diag}$ rises but $\xi$ stays small}

This suggests that the spectrum is moving onto the correct diagonal channel, but not yet condensing into a sharp peak. In simple words, the network may be ``using the right family of modes'' without having developed long-range coherent order.

\subsection{If both $M_{\diag}$ and $\xi$ rise together}

This is stronger evidence that the model is organizing around the true modular-addition structure in a way that resembles a continuous ordering phenomenon.

\subsection{If the peak location wanders during training}

Then the model may still be exploring multiple competing patterns. A stable peak location is often a sign that the model has selected one dominant ordering channel.

\section{How to run a scaling study with different values of $n$}

The user asked specifically how to test scaling relations by training models for different values of $n$. This is the most natural finite-size study here.

\subsection{Why $n$ plays the role of system size}

The pair-state lattice is $\Zn^2$, an $n \times n$ periodic grid. Therefore the natural linear size is
\begin{equation}
    L = n.
\end{equation}
As $n$ grows, the lattice becomes larger, and the momentum spacing $2\pi/n$ becomes finer. This is exactly the type of setting in which finite-size scaling ideas are meaningful.

\subsection{What to keep fixed and what to vary}

A clean scaling study could use values such as
\begin{equation}
    n \in \{8,10,12,16,20,24\}.
\end{equation}
For each $n$, train several random seeds with the same overall recipe:
\begin{itemize}
    \item same optimizer and learning rate schedule,
    \item same training protocol,
    \item same sequence length rule,
    \item same monitoring frequency,
    \item ideally similar effective model capacity relative to task difficulty.
\end{itemize}

In practice, there are two reasonable choices:
\begin{enumerate}
    \item keep the architecture fixed across $n$, or
    \item scale width mildly with $n$ so the model is neither much too strong nor much too weak at large $n$.
\end{enumerate}
The first choice is cleaner for a first study.

\subsection{The main observables to record}

For each $n$ and each training step $t$, record:
\begin{equation}
    M_{\diag}(t;n), \qquad \xi(t;n), \qquad S_t(k,\ell;n).
\end{equation}
Across random seeds, also record the fluctuations of a chosen order parameter. A simple choice is
\begin{equation}
    m(t;n) = \sqrt{S_t(k^*(t);n)},
\end{equation}
where $k^*(t)$ is the dominant nonzero peak.

Then define the susceptibility-like quantity
\begin{equation}
    \chi(t;n) = n^2 \left( \langle m(t;n)^2 \rangle_{\text{seeds}} - \langle m(t;n) \rangle_{\text{seeds}}^2 \right).
\end{equation}
The factor $n^2$ appears because the number of lattice sites is $n^2$.

\subsection{How to locate the pseudocritical point}

For a finite system there is no true singularity, so one usually identifies a \emph{pseudocritical} training step $t_c(n)$ by one of the following:
\begin{itemize}
    \item the step where $\chi(t;n)$ is maximal,
    \item the step where $dM_{\diag}/dt$ is maximal,
    \item the step where $d\xi/dt$ is maximal,
    \item the step where the structure-factor peak becomes clearly dominant.
\end{itemize}
The susceptibility peak is often the cleanest choice.

\subsection{Finite-size scaling forms}

If the training dynamics behaves like a continuous transition, the finite-size scaling ansatz suggests relations of the form
\begin{equation}
    m(t;n) = n^{-\beta/\nu} \, F\!\left((t-t_c)n^{1/\nu}\right),
\end{equation}
\begin{equation}
    \chi(t;n) = n^{\gamma/\nu} \, G\!\left((t-t_c)n^{1/\nu}\right),
\end{equation}
\begin{equation}
    \frac{\xi(t;n)}{n} = H\!\left((t-t_c)n^{1/\nu}\right).
\end{equation}
Here:
\begin{itemize}
    \item $\beta$ describes how the order parameter turns on,
    \item $\gamma$ describes how susceptibility grows,
    \item $\nu$ describes how the correlation length grows.
\end{itemize}
The functions $F$, $G$, and $H$ are unknown scaling functions.

\subsection{A simple crossing test}

A very useful diagnostic is to plot
\begin{equation}
    \xi(t;n)/n
\end{equation}
for many values of $n$ on the same graph. If the system behaves like a continuous transition, these curves often cross near the same training step. That common crossing estimates the critical point.

\subsection{How to estimate $\nu$ numerically}

One practical method is data collapse.

\begin{enumerate}
    \item For each $n$, estimate $t_c(n)$.
    \item Pick a candidate value of $\nu$.
    \item Plot $\xi(t;n)/n$ as a function of
    \begin{equation}
        x = (t-t_c)n^{1/\nu}.
    \end{equation}
    \item Adjust $\nu$ until the curves collapse as well as possible.
\end{enumerate}

Because training step is not an equilibrium control parameter, the resulting exponents should be interpreted carefully. They are best viewed as \emph{effective critical exponents} unless a deeper theoretical mapping justifies a stricter interpretation.

\subsection{How to estimate $\beta$ and $\gamma$}

Once $\nu$ is chosen, you can attempt collapse for the order parameter and susceptibility:
\begin{equation}
    n^{\beta/\nu} m(t;n) \quad \text{vs.} \quad (t-t_c)n^{1/\nu},
\end{equation}
\begin{equation}
    n^{-\gamma/\nu} \chi(t;n) \quad \text{vs.} \quad (t-t_c)n^{1/\nu}.
\end{equation}
If the curves for different $n$ fall approximately on top of each other, that supports the chosen exponents.

\subsection{What if collapse fails?}

That can happen for several reasons:
\begin{itemize}
    \item the monitored observable is not a good order parameter,
    \item the transition is only a crossover,
    \item the range of $n$ is too small,
    \item the model is too overpowered or too underpowered,
    \item the training-step control parameter is not the cleanest one.
\end{itemize}

A very useful next step is then to vary a more conventional control parameter, such as weight decay, train fraction, noise level, or model width, while still monitoring the same spectral quantities.

\section{A compact experimental recipe}

For each $n$:
\begin{enumerate}
    \item Train several random seeds.
    \item Save checkpoints every fixed number of steps.
    \item At each checkpoint, compute the full structure factor $S_t(k,\ell)$.
    \item From $S_t$, compute:
    \begin{itemize}
        \item diagonal spectral mass $M_{\diag}(t)$,
        \item dominant mode $k^*(t)$,
        \item correlation length estimator $\xi(t)$,
        \item order parameter $m(t)=\sqrt{S_t(k^*)}$,
        \item susceptibility-like quantity $\chi(t)$ from seed-to-seed fluctuations.
    \end{itemize}
    \item Plot $M_{\diag}(t)$, $\xi(t)$, and $\chi(t)$.
    \item Search for a common sharp region where these quantities reorganize.
    \item Compare several values of $n$ and test crossing/collapse relations.
\end{enumerate}

\section{Final conceptual summary}

\begin{itemize}
    \item The natural state space for Fibonacci modulo $n$ is the pair lattice $\Zn^2$.
    \item A centered logit table $\widetilde{\ell}(c\mid a,b)$ defines a field on that lattice.
    \item The structure factor is the power spectrum of that field, or equivalently the Fourier transform of its two-point correlation function.
    \item The diagonal spectral mass is the fraction of the structure-factor weight that lies on the diagonal modes $(m,m)$.
    \item Diagonal spectral mass is a useful order-parameter-like summary because the exact modular-addition rule has diagonal Fourier support.
    \item The structure factor contains more information than diagonal spectral mass, especially the width of the dominant peak.
    \item The width of that peak determines a correlation length estimator.
    \item Near a continuous transition, the Ornstein-Zernike form explains why the structure factor near the dominant wavevector looks like
    \begin{equation}
        S(k) \approx \frac{A}{\xi^{-2}+|k-k^*|^2}.
    \end{equation}
    \item Studying the dynamics of $M_{\diag}$, $S(k)$, and $\xi$ during training provides a concrete way to search for critical-like organization in the model.
    \item Repeating the experiment for different $n$ allows finite-size scaling tests and numerical estimates of effective critical exponents.
\end{itemize}

\end{document}
